% 教学内容层；原模板主题、图标及校园素材保持原样。
\usepackage{amsmath,amssymb,booktabs,array,pgfplots}
\pgfplotsset{compat=1.18}
\definecolor{lessonblue}{HTML}{225EA8}
\definecolor{lessonred}{HTML}{BA312D}
\definecolor{lessonpurple}{HTML}{82589F}
\setbeamerfont{normal text}{size=\large}
\setbeamerfont{block title}{size=\normalsize}
\setbeamerfont{block body}{size=\normalsize}
\setbeamertemplate{itemize item}{\color{campusgreen}\small$\bullet$}
\setbeamertemplate{enumerate item}{\color{campusgreen}\insertenumlabel.}
\setlength{\parskip}{.45em}
\setlength{\abovedisplayskip}{.5em}
\setlength{\belowdisplayskip}{.5em}
\newcommand{\R}{\mathbb R}
\newcommand{\Q}{\mathbb Q}
\newcommand{\ans}[1]{{\color{lessonred}#1}}
\newcommand{\key}[1]{{\color{campusgreen}\bfseries #1}}
\newcommand{\reason}[1]{{\color{lessonblue}#1}}
\newcommand{\gap}{\par\vspace{.6em}}
\pgfplotsset{lessonaxis/.style={axis lines=middle,axis line style={->},
 xlabel={$x$},ylabel={$y$},
 tick label style={font=\scriptsize},label style={font=\small},
 width=7.1cm,height=5.1cm,samples=150,clip=true,unbounded coords=jump,
 every axis plot/.append style={thick},
 legend style={font=\scriptsize,draw=none,fill=white,fill opacity=.9,text opacity=1},
 legend cell align=left}}
% 五个基本幂函数，均为可编辑矢量曲线。
\newcommand{\fivepowers}{%
\begin{tikzpicture}\begin{axis}[lessonaxis,xmin=-2.1,xmax=3.2,ymin=-2.1,ymax=3.6,
 xtick={-2,-1,1,2,3},ytick={-2,-1,1,2,3},legend pos=north west]
 \addplot[black,domain=-2:3.1]{x};\addlegendentry{$y=x$}
 \addplot[lessonblue,domain=-1.85:1.85]{x^2};\addlegendentry{$y=x^2$}
 \addplot[lessonred,domain=-1.3:1.5]{x^3};\addlegendentry{$y=x^3$}
 \addplot[campusgreen,domain=0:3.1]{sqrt(x)};\addlegendentry{$y=\sqrt{x}$}
 \addplot[lessonpurple,domain=.285:3.1]{1/x};\addlegendentry{$y=1/x$}
 \addplot[lessonpurple,domain=-2:-.48,forget plot]{1/x};
 \addplot[densely dashed,gray,forget plot] coordinates {(0,1)(1,1)(1,0)};
 \addplot[only marks,mark=*,mark size=1.3pt,forget plot] coordinates {(1,1)};
\end{axis}\end{tikzpicture}}
% 示意值只用于表达 0<q<m<1，未给原题补入数值条件。
\newcommand{\powerorder}{%
\begin{tikzpicture}\begin{axis}[lessonaxis,xmin=0,xmax=3.2,ymin=0,ymax=4,
 xtick={1,2,3},ytick={1,2,3},legend pos=north west]
 \addplot[lessonred,domain=0:1.95]{x^2};\addlegendentry{$y=x^p$}
 \addplot[campusgreen,domain=0:3.1]{x^.6};\addlegendentry{$y=x^m$}
 \addplot[lessonblue,domain=0:3.1]{x^.3};\addlegendentry{$y=x^q$}
 \addplot[lessonpurple,domain=.25:3.1]{1/x};\addlegendentry{$y=x^n$}
 \addplot[densely dashed,gray,forget plot] coordinates {(0,1)(1,1)(1,0)};
\end{axis}\end{tikzpicture}}
\newcommand{\expplot}[1][2]{%
\begin{tikzpicture}\begin{axis}[lessonaxis,xmin=-2.5,xmax=2.5,ymin=0,ymax=5,
 xtick={-2,-1,1,2},ytick={1,2,3,4},legend pos=north west]
 \addplot[lessonred,domain=-2.4:2.3]{2^x};\addlegendentry{$y=2^x$}
 \ifnum#1>1
 \addplot[lessonblue,domain=-2.3:2.4]{2^(-x)};\addlegendentry{$y=(1/2)^x$}
 \fi
 \ifnum#1>2
 \addplot[lessonpurple,domain=-2.4:1.45]{3^x};\addlegendentry{$y=3^x$}
 \addplot[campusgreen,domain=-1.45:2.4]{3^(-x)};\addlegendentry{$y=(1/3)^x$}
 \fi
 \addplot[only marks,mark=*,mark size=1.4pt,black,forget plot] coordinates {(0,1)};
\end{axis}\end{tikzpicture}}
% 四个选项依据原图特征重绘，B 为实立方根。
\newcommand{\rootchoices}{%
\foreach \letter/\kind in {A/1,B/2,C/3,D/4}{%
\begin{minipage}[t]{.24\textwidth}\centering
\begin{tikzpicture}\begin{axis}[lessonaxis,width=3.5cm,height=3.3cm,
 xmin=-1.6,xmax=1.7,ymin=-1.5,ymax=1.8,xtick={1},ytick={1},samples=201]
 \addplot[dashed,gray,domain=-1.4:1.6]{x};
 \ifnum\kind=1 \addplot[black,domain=-1.5:1.5]{tanh(x)};\fi
 \ifnum\kind=2 \addplot[black,domain=-1.5:1.5]{sign(x)*abs(x)^(1/3)};\fi
 \ifnum\kind=3 \addplot[black,domain=-1.12:1.2]{x^3};\fi
 \ifnum\kind=4 \addplot[black,domain=-1.5:1.5]{sign(x+.4)*abs(x+.4)^(1/3)};\fi
\end{axis}\end{tikzpicture}\\[-.5em]\letter
\end{minipage}}\par}
